\begin{abstract}
Before a complex system can adapt or self-organize, it must remain identifiable through time. From biological organisms to computational processes, disparate systems share a single existential mandate: maintaining structural coherence against entropic erasure. We formalize the mechanics of this persistence, deriving the substrate-independent, minimal architecture. We show that any persistent system must satisfy three fundamental constraints: Finiteness (bounded resources), Conservation (preserved information), and Causality (well-defined evolution), yielding six irreducible pillars: Capacity, Map, Protocol, Governor, Toll, and Margin. This static architecture is grounded in time by defining the minimal dynamical description of state transitions: toll ($c$), structural complexity ($K$), and coordinated transition flux ($f$) whose product defines the dissipation density $\mathcal{D} = c \cdot K \cdot f$. Integrating $\mathcal{D}$ over a system's history defines the Operational Action $S_{\Phi}$. The resulting Selection Principle shows that minimum-action configurations asymptotically dominate observable histories. Mapping systems onto this architectural invariant provides a common basis for comparing maintenance constraints, identifying structurally analogous failure modes, and formulating cross-domain hypotheses for subsequent validation.
\end{abstract}